% GENERATED FILE. Edit canonical lessons, then run npm run generate:books.
% canonical-source-hash: fnv1a64:03400f5ae4d21d3e
% canonical-lessons: JA-C98-doushite, JA-C98-riyuu, JA-C98-tame, JA-C98-sei, JA-C98-yue

\chapter{Why, Reason, For the sake of, Fault, Reason}
\label{ch:ja-why-reason-for-the-sake-of-fault-reason}
\hlchaptermodality{98}

\begin{chapteropening}
\textbf{By the end of this chapter:} \emph{I can say why, a reason, for the sake of, the fault and a reason.}

Complete the last lesson of chapter 98: i can say why, a reason, for the sake of, the fault and a reason.
\end{chapteropening}

\section[\texorpdfstring{\ja{どうして} (doushite)}{どうして (doushite)}]{\ja{どうして} (doushite) — why}

\label{lesson:JA-C98-doushite}



\begin{quote}
Before the new one: say the Japanese for a method, then the Japanese for a way of doing.
\end{quote}

\subsection*{You'll want to know: \ja{どうして}}
\textbf{\ja{どうして}} — \emph{doushite} — \textquotedblleft{}why\textquotedblright{}.

\begin{grammarlens}[title={What you've built}]
One more word for this chapter.
\end{grammarlens}

\subsection*{Your turn}
\begin{itemize}
\raggedright
  \item \emph{Say it:} \emph{doushite}
  \item \emph{Say it:} \emph{doushite}, once more
  \item \emph{Say it:} say \emph{doushite}
  \item \emph{Recall:} say the Japanese for a method, then the Japanese for a way of doing, then say \emph{doushite} again
\end{itemize}

\begin{tcolorbox}[breakable,colback=teal!4,colframe=teal!35!black,title={Before you move on}]
What does \ja{どうして} mean? (\textquotedblleft{}Why\textquotedblright{}.) Say it once more.
\end{tcolorbox}
\section[\texorpdfstring{\ja{りゆう} (riyuu)}{りゆう (riyuu)}]{\ja{りゆう} (riyuu) — a reason}

\label{lesson:JA-C98-riyuu}



\begin{quote}
Before the new one: say the Japanese for a way of doing, then the Japanese for why.
\end{quote}

\subsection*{You'll want to know: \ja{りゆう}}
\textbf{\ja{りゆう}} — \emph{riyuu} — \textquotedblleft{}a reason\textquotedblright{}.

\begin{grammarlens}[title={What you've built}]
One more word for this chapter.
\end{grammarlens}

\subsection*{Your turn}
\begin{itemize}
\raggedright
  \item \emph{Say it:} \emph{riyuu}
  \item \emph{Say it:} \emph{riyuu}, once more
  \item \emph{Say it:} say \emph{riyuu}
  \item \emph{Recall:} say the Japanese for a way of doing, then the Japanese for why, then say \emph{riyuu} again
\end{itemize}

\begin{tcolorbox}[breakable,colback=teal!4,colframe=teal!35!black,title={Before you move on}]
What does \ja{りゆう} mean? (\textquotedblleft{}A reason\textquotedblright{}.) Say it once more.
\end{tcolorbox}
\section[\texorpdfstring{\ja{ため} (tame)}{ため (tame)}]{\ja{ため} (tame) — for the sake of; because of}

\label{lesson:JA-C98-tame}



\begin{quote}
Before the new one: say the Japanese for why, then the Japanese for a reason.
\end{quote}

\subsection*{You'll want to know: \ja{ため}}
\textbf{\ja{ため}} — \emph{tame} — \textquotedblleft{}for the sake of; because of\textquotedblright{}.

\begin{grammarlens}[title={What you've built}]
One more word for this chapter.
\end{grammarlens}

\subsection*{Your turn}
\begin{itemize}
\raggedright
  \item \emph{Say it:} \emph{tame}
  \item \emph{Say it:} \emph{tame}, once more
  \item \emph{Say it:} say \emph{tame}
  \item \emph{Recall:} say the Japanese for why, then the Japanese for a reason, then say \emph{tame} again
\end{itemize}

\begin{tcolorbox}[breakable,colback=teal!4,colframe=teal!35!black,title={Before you move on}]
What does \ja{ため} mean? (\textquotedblleft{}For the sake of; because of\textquotedblright{}.) Say it once more.
\end{tcolorbox}
\section[\texorpdfstring{\ja{せい} (sei)}{せい (sei)}]{\ja{せい} (sei) — the fault, the cause}

\label{lesson:JA-C98-sei}



\begin{quote}
Before the new one: say the Japanese for a reason, then the Japanese for for the sake of; because of.
\end{quote}

\subsection*{You'll want to know: \ja{せい}}
\textbf{\ja{せい}} — \emph{sei} — \textquotedblleft{}the fault, the cause\textquotedblright{}.

\begin{grammarlens}[title={What you've built}]
One more word for this chapter.
\end{grammarlens}

\subsection*{Your turn}
\begin{itemize}
\raggedright
  \item \emph{Say it:} \emph{sei}
  \item \emph{Say it:} \emph{sei}, once more
  \item \emph{Say it:} say \emph{sei}
  \item \emph{Recall:} say the Japanese for a reason, then the Japanese for for the sake of; because of, then say \emph{sei} again
\end{itemize}

\begin{tcolorbox}[breakable,colback=teal!4,colframe=teal!35!black,title={Before you move on}]
What does \ja{せい} mean? (\textquotedblleft{}The fault, the cause\textquotedblright{}.) Say it once more.
\end{tcolorbox}
\section[\texorpdfstring{\ja{ゆえ} (yue)}{ゆえ (yue)}]{\ja{ゆえ} (yue) — a reason (formal)}

\label{lesson:JA-C98-yue}



\begin{quote}
Before the new one: say the Japanese for for the sake of; because of, then the Japanese for the fault.
\end{quote}

\subsection*{You'll want to know: \ja{ゆえ}}
\textbf{\ja{ゆえ}} — \emph{yue} — \textquotedblleft{}a reason (formal)\textquotedblright{}.

\begin{grammarlens}[title={What you've built}]
One more word for this chapter.
\end{grammarlens}

\subsection*{Your turn}
\begin{itemize}
\raggedright
  \item \emph{Say it:} \emph{yue}
  \item \emph{Say it:} \emph{yue}, once more
  \item \emph{Say it:} say \emph{yue}
  \item \emph{Recall:} say the Japanese for for the sake of; because of, then the Japanese for the fault, then say \emph{yue} again
\end{itemize}

\begin{tcolorbox}[breakable,colback=teal!4,colframe=teal!35!black,title={Before you move on}]
What does \ja{ゆえ} mean? (\textquotedblleft{}A reason (formal)\textquotedblright{}.) Say it once more.
\end{tcolorbox}
